% !TeX root = ../regression_logistique.tex
\chapter{Optimisation des coefficients}
\label{ch:entrainer}

\orient{%
écrire le gradient et le contrôler contre une mesure par différences finies ;
écrire un pas ; écrire une boucle dont le critère d'arrêt teste une baisse et
rend des coefficients cohérents avec la perte retournée}{%
chapitre~\ref{ch:donnees} ; \annexeref{ann:gradient}}{%
cahier pratique}

Un jeu de coefficients étant évalué, trois étapes séparent ce point d'une boucle
d'entraînement : la direction du déplacement, sa vérification, puis son
itération.

\section{La pente et sa vérification}

La pente du coût par rapport à chaque coefficient s'obtient par une formule
démontrée à l'\annexeref{ann:gradient}. Les trois pentes sortent ensemble d'un
seul passage sur les observations.

Une erreur de signe ou un indice décalé produit ici un programme qui s'exécute,
converge lentement et reste silencieux. La formule se contrôle sans
démonstration, en mesurant ce qu'elle prétend calculer : une pente est un rapport
d'accroissements, obtenu en déplaçant un coefficient dans les deux sens et en
divisant l'écart de coût par la distance parcourue.

\begin{tache}{Gradient et contrôle numérique}
\tlab{Contrat} \texttt{gradient(weights, X, target)} rend la liste des
$\partial J/\partial w_j$, calculée en un seul passage sur \texttt{X}.
\texttt{measured\_slope(weights, X, target, j, h)} rend la même dérivée par
différence centrée, en appelant deux fois \texttt{total\_cost}.

\tlab{Contrôle} En $(0{,}3\,;-0{,}7\,;0{,}4)$ et $h=10^{-5}$ :
\begin{center}\ttfamily\footnotesize
\begin{tabular}{@{}rrr@{}}
colonne & formule & mesure\\
0 & +0.358202569 & +0.358202569\\
1 & +0.071184453 & +0.071184453\\
2 & -0.098350181 & -0.098350181\\
\end{tabular}
\end{center}
Écart attendu : de l'ordre de $10^{-10}$.

\tlab{Indice 1} La formule est
$\frac{1}{m}\sum_i (p_i-y_i)\,x_{ij}$ : accumuler un résidu par ligne, puis le
répartir sur les colonnes de cette ligne.

\tlab{Indice 2} \texttt{measured\_slope} ignore tout de la régression
logistique : elle déplace \texttt{weights[j]} de $\pm h$, rappelle
\texttt{total\_cost}, et divise par $2h$.

\tlab{Transfert} Le point d'évaluation est choisi loin de zéro. Montrer qu'en
$(0,0,0)$ un facteur $1/m$ omis passerait inaperçu, et proposer un second point
de contrôle qui révélerait un indice décalé.
\corrige{step5\_gradient.py}
\end{tache}

\begin{figure}[htbp]
\centering
\begin{graphique}{Un bloc d'entrée se sépare en deux chemins : à gauche la fonction gradient, un passage sur les observations ; à droite measured_slope, deux coûts complets. Les deux convergent vers un bloc de comparaison indiquant que les valeurs coïncident sur neuf décimales.}
\begin{tikzpicture}[
  box/.style={rounded corners=2.5pt,inner sep=6pt,align=center,font=\small,
    text width=62mm,draw=ink!30},
  arr/.style={-{Latex[length=2mm]},thick,ink!60}]
\node[box,fill=gridgray!20] (in)
  {\texttt{weights} $=(0{,}3\,;-0{,}7\,;0{,}4)$\\ \texttt{X}, \texttt{target}};
\node[box,fill=accentlight,below=11mm of in,xshift=-34mm] (form)
  {\texttt{gradient}\\[2pt] la formule,\\ un passage sur les observations};
\node[box,fill=rulelight,below=11mm of in,xshift=34mm] (mes)
  {\texttt{measured\_slope}\\[2pt] deux coûts complets,\\ déplacement de $\pm10^{-5}$};
\coordinate (mid) at ($(form.south)!0.5!(mes.south)$);
\node[box,fill=white,draw=ink!45,below=10mm of mid,text width=100mm] (cmp)
  {les deux valeurs coïncident sur neuf décimales,\\ aux trois colonnes};
\draw[arr] (in) -- (form);
\draw[arr] (in) -- (mes);
\draw[arr] (form) -- (cmp);
\draw[arr] (mes) -- (cmp);
\end{tikzpicture}
\end{graphique}
\caption{Deux chemins indépendants vers la même dérivée. Le point d'évaluation
est choisi loin de zéro, où plusieurs erreurs classiques rendent zéro des deux
côtés.}
\label{fig:controle-gradient}
\end{figure}
\FloatBarrier

L'écart résiduel vient de la mesure. Centrée sur un pas $h$, une différence
commet une erreur en $h^2$, soit $10^{-10}$ pour $h=10^{-5}$, ce que les trois
lignes retrouvent. Un $h$ plus grand mesurerait une corde au lieu d'une tangente ;
un $h$ plus petit ferait perdre à la soustraction ses décimales significatives.

\begin{vigilance}[Ordre de grandeur de l'écart]
Un écart de l'ordre de $10^{-10}$ valide la formule. Un écart de $10^{-2}$
désigne un facteur oublié, un signe opposé désigne une soustraction inversée, et
une seule colonne fausse désigne l'indice.
\end{vigilance}

\section{Un pas}

Descendre consiste à retrancher la pente, multipliée par la longueur du pas.

\begin{tache}{Un pas de descente}
\tlab{Contrat} \texttt{one\_step(weights, X, target)} rend une \emph{nouvelle}
liste de coefficients, chacun diminué de \texttt{ALPHA} fois sa dérivée.

\tlab{Contrôle} À partir de coefficients nuls et $\alpha=1$ :
\begin{center}\ttfamily\footnotesize
\begin{tabular}{@{}llll@{}}
turn 0 & cost 0.693147 & weights & +0.0000 +0.0000 +0.0000\\
turn 1 & cost 0.538507 & weights & -0.2956 -0.2289 +0.1924\\
turn 2 & cost 0.449297 & weights & -0.5194 -0.4005 +0.3364\\
turn 3 & cost 0.393020 & weights & -0.6955 -0.5348 +0.4500\\
\end{tabular}
\end{center}
Le premier biais, $-0{,}2956$, se recalcule à la main : voir ci-dessous.

\tlab{Indice 1} Construire la liste neuve plutôt que modifier en place.
Corriger \texttt{weights[0]} avant de calculer la correction de
\texttt{weights[1]} évaluerait les deux pentes en deux points différents.

\tlab{Indice 2} Le signe est un moins : on se déplace à l'opposé de la pente.
Si le coût monte au premier tour, la cause est le signe ou la valeur d'\texttt{ALPHA}.

\tlab{Transfert} Prédire, sans exécuter, le signe de $w_1$ et de $w_2$ après le
premier tour, à partir de la position des \Gryf dans le nuage
(figure~\ref{fig:nuage}).
\corrige{step6\_one\_step.py}
\end{tache}

Les coefficients étant nuls, le modèle produit $0{,}5$ pour les $1\,600$ lignes,
et la pente du biais vaut la moyenne des résidus $\sigma(z)-y$, où $y$ vaut $1$
pour les \Gryf et $0$ pour les autres.

\begin{figure}[htbp]
\centering
\begin{graphique}{Une barre horizontale partagée en deux : un segment bleu de 327 Gryffindor à gauche, un segment gris de 1273 élèves des autres maisons à droite. Sous chaque segment, le résidu correspondant, -0,5 et +0,5, puis le calcul de leur moyenne pondérée, égale à 0,2956.}
\begin{tikzpicture}[font=\small]
\fill[accent!65] (0,0) rectangle (2.45,0.72);
\fill[gridgray!45] (2.45,0) rectangle (12,0.72);
\draw[ink!35] (0,0) rectangle (12,0.72);
\node[font=\footnotesize\bfseries,text=white] at (1.22,0.36) {$327$};
\node[font=\footnotesize,text=ink] at (7.2,0.36)
  {$1\,273$ élèves des trois autres maisons};
\node[anchor=south,font=\footnotesize,text=accent] at (1.22,0.78)
  {\Gryf, $20{,}44\,\%$};
\node[anchor=north,font=\footnotesize,text=accent] at (1.22,-0.12)
  {$\sigma(0)-1=-0{,}5$};
\node[anchor=north,font=\footnotesize,text=ink] at (7.2,-0.12)
  {$\sigma(0)-0=+0{,}5$};
\node[anchor=north,font=\small] at (6,-1.05)
  {moyenne $=\dfrac{1\,273\times0{,}5-327\times0{,}5}{1\,600}=0{,}2956$};
\end{tikzpicture}
\end{graphique}
\caption{Le pas retranche cette moyenne au biais, d'où $-0{,}2956$ au premier
tour : le biais rejoint la valeur qui fait annoncer $20{,}44\,\%$ à toutes les
observations.}
\label{fig:premierpas}
\end{figure}
\FloatBarrier

\section{Itération jusqu'à stagnation}

Un pas isolé fait baisser la perte d'une fraction. La descente répète ce pas
jusqu'à ce que la baisse devienne négligeable, ce qui demande une boucle et deux
conditions d'arrêt.

\begin{tache}{La boucle et ses deux arrêts}
\tlab{Contrat} \texttt{descend(X, target)} rend le couple
\texttt{(weights, cost)} correspondant au même état, ainsi que le nombre de tours
consommés. La boucle s'arrête sur une baisse relative inférieure à
\texttt{EPSILON} $=10^{-6}$, ou sur un plafond de $6\,000$ tours.

\tlab{Contrôle}
\begin{center}\ttfamily\footnotesize
stopped after 983 iterations, cost 0.107257\\
weights: -5.1755 -3.7460 +3.7874
\end{center}
Le nombre d'itérations se lit en premier : à $983$ pour un plafond de $6\,000$,
c'est le critère de stagnation qui a mis fin à la boucle. Une valeur égale au
plafond signalerait un pas trop court, ou l'absence de minimum fini
(\annexeref{ann:norme}).

\tlab{Indice 1} Le test doit porter sur \texttt{previous - cost}, une baisse, et
non sur \texttt{abs(previous - cost)}.

\tlab{Indice 2} Appliquer le pas \emph{avant} de calculer la perte retournée :
les deux valeurs décrivent alors le même état.

\tlab{Transfert} Le seuil relatif $10^{-6}$ est atteint après $983$ tours, alors
que la frontière a cessé de se déplacer bien avant (ch.~\ref{ch:descente}).
Proposer un critère d'arrêt qui testerait la stabilité des \emph{décisions}, et
dire ce qu'il coûterait par tour.
\corrige{step7\_loop.py}
\end{tache}

\begin{vigilance}[Trois pièges du critère d'arrêt]
\textbf{Valeur absolue.} Écrire \texttt{abs(previous - cost) < seuil} accepte
aussi une hausse faible de la perte et peut donc arrêter à tort une boucle
instable.

\textbf{Relatif ou absolu.} Écrire \texttt{... / max(1, abs(cost))} ne
normalise pas le coût lorsqu'il est inférieur à $1$. La forme
\texttt{previous - cost < EPSILON * abs(previous)} conserve une interprétation
relative.

\textbf{Cohérence entre retours.} La perte et les coefficients retournés
doivent décrire le même état ; appliquer un pas après le dernier calcul de perte
les décale d'une itération.
\end{vigilance}

\begin{figure}[htbp]
\centering
\begin{graphique}{Trois blocs empilés reliés en boucle : total_cost, gradient, one_step, la sortie de one_step revenant sur total_cost avec la mention au plus 6000 tours. Une flèche latérale sort de total_cost vers le retour des coefficients et du coût, sous condition d'une baisse relative inférieure à 10 puissance -6.}
\begin{tikzpicture}[
  st/.style={draw=ink!30,rounded corners=2.5pt,align=center,font=\small,
    inner sep=6pt,text width=34mm},
  lab/.style={font=\footnotesize,text=ink,align=left},
  arr/.style={-{Latex[length=2mm]},thick,ink!60}]
\node[st,fill=accentlight] (c) {\texttt{total\_cost}};
\node[st,fill=rulelight,below=13mm of c] (g) {\texttt{gradient}};
\node[st,fill=warninglight,below=13mm of g] (u) {\texttt{one\_step}};
\node[st,fill=white,draw=accent!60,left=16mm of c] (out)
  {\texttt{return}\\ \texttt{weights, cost}};
\coordinate (ret) at ($(u.east)+(12mm,0)$);
\draw[arr] (c) -- (g);
\draw[arr] (g) -- (u);
\draw[arr] (u.east) -- (ret) |- (c.east);
\node[lab,right=1mm of ret,text width=26mm] {au plus\\ $6\,000$ tours};
\draw[arr] (c) -- node[lab,above=1mm,align=center]
  {baisse relative\\ sous $10^{-6}$} (out);
\end{tikzpicture}
\end{graphique}
\caption{Le coût est calculé une fois par tour, pour le seul test d'arrêt. Le
gradient constitue un second passage sur les $1\,600$ lignes : une itération en
coûte deux.}
\label{fig:boucle}
\end{figure}
\FloatBarrier

Le seuil relatif arrête les itérations lorsque la baisse du coût devient
négligeable. Le plafond couvre les configurations où les coefficients divergent
et où ce seuil n'est jamais atteint (\annexeref{ann:norme}).

Ces coefficients constituent l'état \Mstop. Les figures de la partie~I emploient
\Mquatre, arrêté au quatre-centième tour : $-4{,}745$, $-3{,}454$, $+3{,}469$.
Les deux états diffèrent d'environ $9\,\%$ en norme, et leurs rapports coïncident
à trois millièmes près : la frontière se stabilise avant la perte, tandis que
les derniers tours resserrent la sigmoïde.
